\chapter{Random Variables for Econometrics}\label{ch:b1:random-variables-for-econometrics}

Draw five men at random from the 2,490 of the PSID sample and average their
1978 earnings; draw five more and average again. Repeated 10,000 times, the
averages spread with a standard deviation of 6,979 dollars around the
population mean of 21,554 (dollars of 1982, as everywhere in this chapter's
earnings). % ledger: b1:nsw-1982dollars
With thirty men the spread falls to 2,840; with two
hundred, to 1,093. % figdata: means_stats
% figdata: pop
Each figure is close to the population's standard deviation divided by
the square root of the sample size. This chapter gives the probability
behind that rule --- \omterm{def:b1:random-variables-for-econometrics:conditional-expectation}{conditional expectations}, the distributions that
tests use, the law of large numbers and the central limit theorem --- and
the cases in which it fails.

\begin{openingchart}{Sample means of 1978 earnings, 10,000 samples of 5 (light),
30 (blue) and 200 (red) men drawn with replacement from the PSID sample,
densities per 10,000 dollars; the population mean is dashed.}{National
Supported Work comparison sample (PSID); simulation with a fixed seed}
\begin{tikzpicture}
\begin{axis}[omrecord, xlabel={Sample mean of earnings (dollars of 1982)}, ylabel={Density}, xmin=0, xmax=50000, ymin=0,
  ymax=3.8, xtick={0,10000,...,50000}, width=0.7\linewidth, height=0.38\linewidth, scaled x ticks=false]
  \addplot[gray, thick, const plot mark mid] table[x=mid, y=n5, col sep=comma] {figdata/bachelor-1/30-random-variables-for-econometrics/out/means.csv};
  \addplot[omDef, thick, const plot mark mid] table[x=mid, y=n30, col sep=comma] {figdata/bachelor-1/30-random-variables-for-econometrics/out/means.csv};
  \addplot[omThm, thick, const plot mark mid] table[x=mid, y=n200, col sep=comma] {figdata/bachelor-1/30-random-variables-for-econometrics/out/means.csv};
  \draw[black, dashed] (axis cs:21553.9,0) -- (axis cs:21553.9,3.8);
\end{axis}
\end{tikzpicture}
\end{openingchart}

The earnings themselves are skewed to the right (\cref{ch:b1:describing-data}); the
means of five are less skewed, and the means of two hundred look like a bell
curve centred on the population mean.

\begin{recall}
Random variables, expectation, variance, independence, the binomial and
normal laws and the law of large numbers were met in the mathematics
series at school; the central limit theorem is proved in its third
university volume. Describing data is \cref{ch:b1:describing-data}.
\end{recall}

\section{Expectation, variance, covariance}\label{sec:b1:random-variables-for-econometrics:moments}

For random variables $X$, $Y$ and numbers $a$, $b$: $\E[aX + b] = a\E[X] + b$, $\E[X + Y] = \E[X] + \E[Y]$
and $\Var(aX + b) = a^2\Var(X)$.

\begin{notation}
The covariance of two random variables is $\Cov(X, Y) = \E[(X - \E X)(Y - \E Y)]$, defined in
the mathematics series (second university volume). Then
$\Var(X + Y) = \Var(X) + \Var(Y) + 2\Cov(X, Y)$, and for independent $X_1, \dots, X_n$ with variance $\sigma^2$,
$\Var(\bar X) = \sigma^2/n$.
\end{notation}

\section{Joint and conditional distributions}\label{sec:b1:random-variables-for-econometrics:joint}

\begin{definition}[Joint and conditional distributions]\label{def:b1:random-variables-for-econometrics:joint}
\index{joint distribution}\index{conditional distribution}
The \emph{joint distribution} of two discrete random variables gives
$P(X = x, Y = y)$ for every pair; the marginal distributions are its row and
column sums. The \emph{conditional distribution} of $Y$ given $X = x$ is
$P(Y = y \mid X = x) = P(X = x, Y = y)/P(X = x)$.
\end{definition}

\omcode{code/bachelor-1/30-random-variables-for-econometrics/listings/joint_table.py}{1}{14}{A \omterm{def:b1:random-variables-for-econometrics:joint}{joint
distribution} from a sample, with its margins, and a \omterm{def:b1:random-variables-for-econometrics:joint}{conditional distribution}.}\label{lst:b1:random-variables-for-econometrics:joint}

\omoutput{code/bachelor-1/30-random-variables-for-econometrics/outputs/joint_table.txt}{Of men with more
than 12 years of schooling, $39.2\%$ earned over 30,000 dollars; of men with
fewer than 12, $10.3\%$.}\label{out:b1:random-variables-for-econometrics:joint}

\section{Conditional expectation}\label{sec:b1:random-variables-for-econometrics:cef}

\begin{definition}[Conditional expectation]\label{def:b1:random-variables-for-econometrics:conditional-expectation}
\index{conditional expectation}
The \emph{conditional expectation} of $Y$ given $X = x$ is $\E[Y \mid X = x] = \sum_y y\,P(Y = y \mid X = x)$.
As a function of $x$ it is the conditional expectation function; $\E[Y \mid X]$ is that function
applied to $X$, a random variable.
\end{definition}

\begin{theorem}[Law of iterated expectations]\label{thm:b1:random-variables-for-econometrics:lie}
\index{law of iterated expectations}
$\E\big[\E[Y \mid X]\big] = \E[Y]$.
\end{theorem}

\begin{proof}
In the discrete case,
$\sum_x P(X = x)\sum_y y\,P(Y = y \mid X = x) = \sum_x\sum_y y\,P(X = x, Y = y) = \sum_y y\,P(Y = y) = \E[Y]$.
\end{proof}

\omcode{code/bachelor-1/30-random-variables-for-econometrics/listings/iterated.py}{1}{9}{The law
checked on a sample: the mean of the \omterm{def:b1:describing-data:conditional-mean}{conditional means}, weighted by group
shares.}\label{lst:b1:random-variables-for-econometrics:iterated}

\omoutput{code/bachelor-1/30-random-variables-for-econometrics/outputs/iterated.txt}{Equal to the last
printed digit.}\label{out:b1:random-variables-for-econometrics:iterated}

\begin{proposition}[The best predictor]\label{prop:b1:random-variables-for-econometrics:best}
Among all functions $g$, $\E[(Y - g(X))^2]$ is smallest for $g(X) = \E[Y \mid X]$.
\end{proposition}

\begin{proof}
Write $m(X) = \E[Y \mid X]$. Then $\E[(Y - g)^2] = \E[(Y - m)^2] + \E[(m - g)^2] + 2\E[(Y - m)(m - g)]$. Conditionally on
$X = x$, $m - g$ is a number and $\E[Y - m \mid X = x] = 0$, so the last term is zero by the law of
iterated expectations. The middle term is at least zero, and zero for $g = m$.
\end{proof}

\omcode{code/bachelor-1/30-random-variables-for-econometrics/listings/best_predictor.py}{1}{14}{Three
predictors of earnings from schooling.}\label{lst:b1:random-variables-for-econometrics:best}

\omoutput{code/bachelor-1/30-random-variables-for-econometrics/outputs/best_predictor.txt}{The \omterm{def:b1:describing-data:conditional-mean}{conditional
mean} predicts best in the sample; a line comes close.}\label{out:b1:random-variables-for-econometrics:best}

\begin{omfigure}
\begin{tikzpicture}
\begin{axis}[omscatter, xlabel={Years of schooling}, ylabel={Mean earnings (dollars of 1982)}, xmin=7.5, xmax=17.5,
  ymin=10000, ymax=35000, width=0.56\linewidth, height=0.36\linewidth, xtick={8,...,17},
  scaled y ticks=false]
  \addplot[only marks, mark=*, mark size=2pt, omDef] table[x=education, y=re78, col sep=comma] {figdata/bachelor-1/30-random-variables-for-econometrics/out/cef.csv};
  \addplot[omThm, thick, sharp plot, mark=none] table[x=education, y=line, col sep=comma] {figdata/bachelor-1/30-random-variables-for-econometrics/out/cef.csv};
\end{axis}
\end{tikzpicture}

{\small The \omterm{def:b1:describing-data:conditional-mean}{conditional means} of earnings by schooling (dots) and the
least-squares line (red): the line approximates the \omterm{def:b1:random-variables-for-econometrics:conditional-expectation}{conditional
expectation} function, with a mean squared error only $1.4\%$ larger.
% computed: code/bachelor-1/30-random-variables-for-econometrics/outputs/best_predictor.txt
}
\end{omfigure}

\begin{definition}[Mean independence]\label{def:b1:random-variables-for-econometrics:mean-independence}
\index{mean independence}
$Y$ is \emph{mean independent} of $X$ when $\E[Y \mid X = x]$ does not depend on $x$. Independence
implies mean independence, which implies $\Cov(X, Y) = 0$; neither converse holds
(\cref{exo:b1:random-variables-for-econometrics:mean-indep}).
\end{definition}

\section{The tools}\label{sec:b1:random-variables-for-econometrics:laws}

\begin{definition}[Chi-square, Student's $t$ and $F$ distributions]\label{def:b1:random-variables-for-econometrics:laws}
\index{chi-square distribution}\index{Student's t distribution}\index{F distribution}
Let $Z, Z_1, \dots, Z_k$ be independent standard normal variables. $V = Z_1^2 + \dots + Z_k^2$ has
the \emph{chi-square distribution} with $k$ degrees of freedom, $\chi^2_k$. If $V$ is independent
of $Z$, $Z/\sqrt{V/k}$ has \emph{Student's $t$ distribution} with $k$ degrees of freedom. If
$V_1 \sim \chi^2_k$ and $V_2 \sim \chi^2_m$ are independent, $(V_1/k)/(V_2/m)$ has the \emph{$F$ distribution} with $k$
and $m$ degrees of freedom.
\end{definition}

\begin{omfigure}
\begin{tikzpicture}
\begin{axis}[omrecord, grid=none, xlabel={$x$}, ylabel={Density}, xmin=-4, xmax=4, ymin=0, ymax=0.42,
  width=0.6\linewidth, height=0.36\linewidth, samples=121, domain=-4:4,
  legend style={font=\scriptsize, at={(0.98,0.98)}, anchor=north east}]
  \addplot[black, thick] {exp(-x^2/2)/sqrt(2*pi)};
  \addplot[omThm, thick, dashed] {1/(2*sqrt(2))*(1 + x^2/2)^(-1.5)};
  \addplot[omDef, thick] {0.37960669*(1 + x^2/5)^(-3)};
  \addplot[gray, thick, dotted] {0.39563218*(1 + x^2/30)^(-15.5)};
  \legend{normal, $t_2$, $t_5$, $t_{30}$}
\end{axis}
\end{tikzpicture}

{\small Student's $t$ against the normal: with 2 degrees of freedom the tails are
much heavier; with 30 the two curves are hard to tell apart.}
\end{omfigure}

The constants are the exact normalising factors: $1/(2\sqrt2)$ for $t_2$, and
$\Gamma(\frac{k+1}2)/(\sqrt{k\pi}\,\Gamma(\frac k2))$, which is $0.3796$ for $k = 5$ and $0.3956$ for $k = 30$.

\omcode{code/bachelor-1/30-random-variables-for-econometrics/listings/distributions.py}{1}{16}{The
critical values of the next two chapters.}\label{lst:b1:random-variables-for-econometrics:critical}

\omoutput{code/bachelor-1/30-random-variables-for-econometrics/outputs/distributions.txt}{Upper-tail
\omterm{def:b1:describing-data:quantile}{quantiles} at $5\%$ and $2.5\%$.}\label{out:b1:random-variables-for-econometrics:critical}

\omcode{code/bachelor-1/30-random-variables-for-econometrics/listings/build_t.py}{1}{14}{Student's $t$ built
from its definition.}\label{lst:b1:random-variables-for-econometrics:build-t}

\omoutput{code/bachelor-1/30-random-variables-for-econometrics/outputs/build_t.txt}{The simulated
\omterm{def:b1:describing-data:quantile}{quantiles} match the exact ones to within $0.03$.}\label{out:b1:random-variables-for-econometrics:build-t}

\begin{measure}[Drawing a sample in a computer]
A computer draws ``random'' numbers from a deterministic sequence that
starts from a \emph{seed}. The same seed gives the same draws, so a
simulation can be rerun and give the same digits, which is why every
simulation in this book fixes its seed. A different seed gives different
digits and, if the simulation is large enough, the same conclusions.
\end{measure}

\section{The law of large numbers and the central limit theorem}\label{sec:b1:random-variables-for-econometrics:clt}

\begin{definition}[Sampling distribution]\label{def:b1:random-variables-for-econometrics:sampling-distribution}
\index{sampling distribution}
The \emph{sampling distribution} of a statistic is its distribution over
repeated samples of the same size from the same population.
\end{definition}

By the law of large numbers, the mean of $n$ independent draws converges to
the population mean as $n$ grows.

\omcode{code/bachelor-1/30-random-variables-for-econometrics/listings/lln.py}{1}{12}{A running mean
of random draws.}\label{lst:b1:random-variables-for-econometrics:lln}

\omoutput{code/bachelor-1/30-random-variables-for-econometrics/outputs/lln.txt}{After 100,000 draws the
mean is within about 30 dollars of the population's.}\label{out:b1:random-variables-for-econometrics:lln}

\begin{theorem}[Central limit theorem]\label{thm:b1:random-variables-for-econometrics:clt}
If $X_1, X_2, \dots$ are independent with the same distribution, mean $\mu$ and finite variance
$\sigma^2$, then for every $z$, $P\big(\sqrt n(\bar X_n - \mu)/\sigma \le z\big) \to \Phi(z)$ as $n \to \infty$, where $\Phi$ is the standard
normal distribution function.
\end{theorem}

\admitted{Proved in the mathematics series, third university volume.}

\omcode{code/bachelor-1/30-random-variables-for-econometrics/listings/clt.py}{1}{17}{The opening chart's
numbers: the spread and the \omterm{def:b1:describing-data:skewness}{skewness} of the sample means.}\label{lst:b1:random-variables-for-econometrics:clt}

\omoutput{code/bachelor-1/30-random-variables-for-econometrics/outputs/clt.txt}{The spread follows
$\sigma/\sqrt n$; the \omterm{def:b1:describing-data:skewness}{skewness} falls from $1.24$ to $0.12$.}\label{out:b1:random-variables-for-econometrics:clt}

\begin{definition}[Random sample]\label{def:b1:random-variables-for-econometrics:random-sample}
\index{random sampling}\index{random sample}
A \emph{random sample} of size $n$ is a set of $n$ observations drawn
independently from one population, each with the population's
distribution; drawing one is \emph{random sampling}.
\end{definition}

\begin{model}{Random sampling}
\assumes Observations drawn independently from one population with a
finite variance.
\predicts Sample means centred on the population mean, with standard
deviation $\sigma/\sqrt n$, and approximately normal for large $n$.
\tested Checked, not tested, by simulation from a real, skewed
population (opening chart): the standard deviations of the means are
within $1\%$ of $\sigma/\sqrt n$ for 5, 30 and 200 men, % figdata: means_stats
but the draws are independent by construction. Whether a real survey's
observations are independent is a property of how it was drawn; a test
compares its estimates with a known total, or the spread of its estimates
with the spread its design implies (\cref{ch:b1:sampling-and-inference}).
\breaks Dependent draws --- months of one \omterm{def:b1:economic-data-python:shapes}{time series}, members of one
family, pupils of one class --- make the spread larger than $\sigma/\sqrt n$ (the
second university year treats them); populations so heavy-tailed that the
variance does not exist make the means settle far more slowly.
\end{model}

\omcode{code/bachelor-1/30-random-variables-for-econometrics/listings/heavy_tails.py}{1}{10}{Twenty
sample means of a population without a variance.}\label{lst:b1:random-variables-for-econometrics:heavy}

\omoutput{code/bachelor-1/30-random-variables-for-econometrics/outputs/heavy_tails.txt}{Ten thousand
times more data narrow the range of the means only by a factor of
about 36.}\label{out:b1:random-variables-for-econometrics:heavy}

With a finite variance, ten thousand times more data would narrow the
range a hundredfold. Here, from samples of 100 to samples of a million, the
range shrinks from $10.12$ to $0.28$, a factor of 36. % computed: code/bachelor-1/30-random-variables-for-econometrics/outputs/heavy_tails.txt

\section*{Exercises}\label{sec:b1:random-variables-for-econometrics:exercises}

\begin{exercise}[$\star$]\label{exo:b1:random-variables-for-econometrics:linear}
$\E[X] = 5$ and $\Var(X) = 2$. Find $\E[2X + 3]$ and $\Var(2X + 3)$.
\end{exercise}

\begin{exercise}[$\star$]\label{exo:b1:random-variables-for-econometrics:cond}
From the rounded joint table of \cref{out:b1:random-variables-for-econometrics:joint}, compute
$P(\text{over } 30{,}000 \mid \text{more than 12 years})$.
\end{exercise}

\begin{exercise}[$\star$]\label{exo:b1:random-variables-for-econometrics:read}
From \cref{out:b1:random-variables-for-econometrics:clt}, give the standard deviation of the
means for $n = 30$ and compare it with $\sigma/\sqrt{30}$.
\end{exercise}

\begin{exercise}[$\star$]\label{exo:b1:random-variables-for-econometrics:predict}
Predict the standard deviation of the means for $n = 120$.
\end{exercise}

\begin{exercise}[$\star\star$]\label{exo:b1:random-variables-for-econometrics:two-groups}
$40\%$ of a population is in group A, with mean earnings 10, and $60\%$ in group B,
with mean 20. Find mean earnings.
\end{exercise}

\begin{exercise}[$\star\star$]\label{exo:b1:random-variables-for-econometrics:var-sum}
$\Var(X) = 4$, $\Var(Y) = 9$, $\Cov(X, Y) = 1.5$. Find $\Var(X + Y)$.
\end{exercise}

\begin{exercise}[$\star\star$]\label{exo:b1:random-variables-for-econometrics:interval}
Using the central limit theorem and the population values of
\cref{out:b1:random-variables-for-econometrics:clt}, give the interval that contains $95\%$ of
the means of 200 men.
\end{exercise}

\begin{exercise}[$\star\star$]\label{exo:b1:random-variables-for-econometrics:chi2}
Find the mean and the variance of $\chi^2_k$, using $\E[Z^2] = 1$ and $\E[Z^4] = 3$.
\end{exercise}

\begin{exercise}[$\star\star$]\label{exo:b1:random-variables-for-econometrics:brothers}
A survey of $n$ men samples $n/2$ pairs of brothers, whose earnings have a
correlation $\rho > 0$ within each pair. Predict the variance of the mean.
\end{exercise}

\begin{exercise}[$\star\star\star$]\label{exo:b1:random-variables-for-econometrics:best}
Prove \cref{prop:b1:random-variables-for-econometrics:best} by writing $\E[(Y - c)^2 \mid X = x]$ as a
function of the number $c$, for each $x$.
\end{exercise}

\begin{exercise}[$\star\star\star$]\label{exo:b1:random-variables-for-econometrics:mean-indep}
Show that \omterm{def:b1:random-variables-for-econometrics:mean-independence}{mean independence} implies zero covariance, and give an example
of zero covariance without \omterm{def:b1:random-variables-for-econometrics:mean-independence}{mean independence}.
\end{exercise}

\begin{exercise}[$\star\star\star$]\label{exo:b1:random-variables-for-econometrics:fail}
What simulation result would show that the normal approximation fails for
means of 30 draws from a given population?
\end{exercise}

\section{Problem: The Brenland Lottery of Surveys}\label{sec:b1:random-variables-for-econometrics:problem}

\begin{problem}[{Weekend problem --- a population, and 10,000 surveys of it}]\label{pb:b1:random-variables-for-econometrics:lottery}
Treat the 2,000 households of Brenland's survey (invented) as the whole
population. Survey firms draw 100 households with replacement; a script
simulates 10,000 such surveys with a fixed seed.

\medskip\noindent\textbf{Part I --- A joint distribution.} Classify households by
size (1 to 5) and by \omterm{def:b1:national-accounts-in-depth:gni}{disposable income} (under 20,000, 20,000 to 30,000, 30,000 or
more dalens).
\begin{enumerate}
  \item Give the marginal distribution of size.
  \item Give $P(\text{size } 1, \text{ income } 30{,}000 \text{ or more})$.
  \item Give $P(\text{income } 30{,}000 \text{ or more} \mid \text{size } 1)$ and $\mid \text{size } 4$.
  \item Are size and income independent?
  \item Why might larger households have higher income?
\end{enumerate}

\medskip\noindent\textbf{Part II --- Conditional expectations.}
\begin{enumerate}[resume]
  \item Give mean income by household size.
  \item Check the law of iterated expectations.
  \item Is income mean independent of size?
  \item Which predictor of income from size has the smallest mean squared error?
  \item Would mean income per person behave the same way?
\end{enumerate}

\medskip\noindent\textbf{Part III --- The sampling distribution.}
\begin{enumerate}[resume]
  \item Give the population mean and standard deviation.
  \item Predict the standard deviation of the survey means.
  \item Give the simulated standard deviation.
  \item Within what distance of the true mean do $95\%$ of survey means fall, by simulation?
  \item And by the central limit theorem?
  \item How would the distance change with 400 households?
  \item Why draw with replacement?
\end{enumerate}

\medskip\noindent\textbf{Part IV --- A $t$ statistic.} In each survey, compute
$t = (\bar y - \mu)/(s/\sqrt{100})$ with $s$ the survey's standard deviation.
\begin{enumerate}[resume]
  \item What distribution would $t$ have if incomes were normal?
  \item Give the $2.5\%$ upper \omterm{def:b1:describing-data:quantile}{quantile} of $t_{99}$.
  \item What share of the surveys has $|t|$ above it?
  \item Give the standard deviation of the 10,000 values of $t$.
  \item Why is the share close to $5\%$ although incomes are skewed?
  \item What would make it far from $5\%$?
  \item What does a survey firm learn from one survey's $t$, when $\mu$ is unknown?
  \item State the result: within what distance of the truth is the mean income of 100 households in $95\%$ of surveys?
\end{enumerate}
\end{problem}
